\documentclass[../../../hsm_tok_q.tex]{subfiles}

\begin{document}
    \setcounter{figure}{0}
    \textsf{図\ref{fig/hsm_tok_2025_2nd_03_01_q}}で，点Oは原点，%
    曲線$\ell$は関数 $y=\myfraca{1}{2}x^2$ のグラフを表している．
    
    点Aは曲線$\ell$上にあり，$x$座標は$-6$である．
        
    曲線$\ell$上にある点をPとする．
    
    次の各問に答えよ．
    \\
    \begin{minipage}{\linewidth}
        \vspace{.5\baselineskip}
        {\centering
        \setlength{\abovecaptionskip}{5pt}
        \includegraphics%
            {\subfix{fig_hsm_tok_2025_2nd_03/fig_hsm_tok_2025_2nd_03_01_q.pdf}}
            \captionof{figure}{} \label{fig/hsm_tok_2025_2nd_03_01_q}
        }
    \end{minipage}
    \vspace{.1\baselineskip}
    %
    \begin{enumerate}
        \item 次の\:\myboxa{①}\:と\:\myboxa{②}\:に当てはまる数を，%
        下の\textsf{ア}～\textsf{ク}のうちからそれぞれ選び，記号で答えよ．

        点Pの$x$座標を$a$，$y$座標を$b$とする．
        
        $a$のとる値の範囲が $-4\jleqq a\jleqq 1$ のとき，$b$のとる値の範囲は，%
        $\:\myboxa{①}\:\jleqq b\jleqq\:\myboxa{②}$ である．
        
        {\renewcommand{\arraystretch}{1.6}%
            \vspace{-.2\baselineskip}
            \begin{tabularx}{\dimexpr\linewidth-\zw}{@{}LLLL@{}}
                {\textsf ア}\; $-16$ & {\textsf イ}\; $-8$&
                {\textsf ウ}\; $-2$ & {\textsf エ}\; $0$ \\
                {\textsf オ}\; $\myfraca{1}{2}$ & {\textsf カ}\; $2$&
                {\textsf キ}\; $8$ & {\textsf ク}\; $16$
            \end{tabularx}
        }

        \item 次の\:\myboxa{③}\:と\:\myboxa{④}\:に当てはまる数を，%
        下の\textsf{ア}～\textsf{エ}のうちからそれぞれ選び，記号で答えよ．

        点Pの$x$座標が$2$のとき，2点A，Pを通る直線の式は，%
        \:$y=\myboxa{③}\:x+\:\myboxa{④}$ である．
        
        {\renewcommand{\arraystretch}{1.6}%
            \vspace{-.2\baselineskip}
                \begin{tabularx}{\dimexpr\linewidth-\zw}{@{}cLLLL@{}}
                \myboxa{③} & {\textsf ア}\; $-2$ & {\textsf イ}\; $-\myfraca{1}{2}$&
                {\textsf ウ}\; $\myfraca{1}{2}$ & {\textsf エ}\; $2$ \\
                \myboxa{④} & {\textsf ア}\; $3$ & {\textsf イ}\; $6$&
                {\textsf ウ}\; $15$ & {\textsf エ}\; $21$
            \end{tabularx}
        }
        
        \newpage

        \item \textsf{図\ref{fig/hsm_tok_2025_2nd_03_02_q}}は，%
            \textsf{図\ref{fig/hsm_tok_2025_2nd_03_01_q}}において，%
            点Pの$x$座標が6より小さい正の数のとき，%
            点Aを通り$x$軸に平行な直線を引き，$y$軸との交点をB，%
            点Aを通り$y$軸に平行な直線を引き，$x$軸との交点をCとし，%
            点Aと点P，点Bと点P，点Cと点Pをそれぞれ結んだ場合を表している．

            $\triangle\mathrm{ACP}$の面積が，$\triangle\mathrm{APB}$の面積の6倍となるとき，%
            点Pの$x$座標を求めよ．
            \\
            \begin{minipage}{\linewidth}
                \vspace{.5\baselineskip}
                {\centering
                    \setlength{\abovecaptionskip}{5pt}
                    \includegraphics%
                        {\subfix{fig_hsm_tok_2025_2nd_03/fig_hsm_tok_2025_2nd_03_02_q.pdf}}
                        \captionof{figure}{} \label{fig/hsm_tok_2025_2nd_03_02_q}
                }
            \end{minipage}
    \end{enumerate}
\end{document}

